In the language of group theory (see \S~96), the general condition that determines whether terms can cross is that they must belong to different irreducible representations of the symmetry group of the system's Hamiltonian\footnote{The electronic terms of the $\mathrm H_2^+$ ion appear to be an exception to this rule. These terms are specified by the angular-momentum projection $\Lambda$ and two elliptic quantum numbers $n_\xi$, $n_\eta$ (see the problem in \S~78). Since all these numbers are associated with functions of different variables, there is in general nothing to prevent crossings between terms $E(R)$ that have the same $\Lambda$ but different pairs $n_\xi$, $n_\eta$, even though their symmetry under rotations and reflections is the same. In fact, however, the separability of variables in the Schrödinger equation for this system means that its Hamiltonian possesses a higher symmetry than its geometrical properties alone imply; relative to this full symmetry group, states with different values of $n_\xi$, $n_\eta$ belong to different types.}.
